\usepackage{catchfile}
\usepackage{ifthen}
\CatchFileDef\TheLanguage{../language.tex}{\endlinechar=-1}
\ifthenelse{\equal{\TheLanguage}{french}}
 {
  \usepackage[french]{babel}
 }
 {
  \usepackage[english]{babel}	
 }
\usepackage{iflang}
\usepackage[utf8]{inputenc}
\usepackage[T1]{fontenc}
\usepackage{csquotes}
\usepackage{lmodern}
\usepackage[
	natbib=true,
	backend=biber,
	style=numeric,
	sorting=none,
  backref=true,
	sortcites
]{biblatex}
\addbibresource{../bibliographie.bib}
\usepackage[a4paper,left=3cm,right=2cm,top=2.5cm,bottom=2.5cm, twoside]{geometry}
\usepackage{xcolor}
%%%%%% MATHS
\usepackage{stmaryrd}
\usepackage{amsmath}
\usepackage{amsfonts}
\usepackage{amssymb}
\usepackage{amsthm}
\usepackage{latexsym}
\usepackage{eurosym}
%%%%%% FONT
\usepackage{libertine}
		%other font
		%\usepackage{euscript}
		%\usepackage{avant}
		%\usepackage{mathptmx}
		%\usepackage{microtype}
		%\usepackage{bm}
		%\usepackage{bbm} 
\usepackage[pdftex]{graphicx}
\usepackage{caption}
\usepackage{subcaption}
%\usepackage{nicematrix}
\usepackage{float}
\usepackage{multirow}
\usepackage{apalike}
%\usepackage{minitoc}
\usepackage{url}
%\usepackage[normalem]{ulem}
\usepackage{tikz}
\usepackage{enumitem}
\usepackage{epigraph}
\usepackage{lscape}
\usepackage{titletoc}
\usepackage{booktabs}
\usepackage{diagbox}
\usepackage{arydshln}
\usepackage{calc}
\usepackage[]{titlesec} 
\usepackage{fancyhdr}
\pagestyle{fancy}
\fancyhead[C]{\rule{.5\textwidth}{4\baselineskip}}% Add something BIG in the header
\setlength{\headheight}{14pt}
\usepackage{lipsum}
\usepackage[]{todonotes}
\usepackage{enumitem} % Customize lists
\usepackage[framemethod=TikZ]{mdframed}
\usepackage{listings}
\usepackage[lined,ruled,onelanguage,commentsnumbered,linesnumbered,inoutnumbered]{algorithm2e}
\usepackage{setspace}\onehalfspacing
\AtBeginDocument{%
  \addtolength\abovedisplayskip{-0.5\baselineskip}%
  \addtolength\belowdisplayskip{-0.5\baselineskip}%
  \addtolength\abovedisplayshortskip{-0.5\baselineskip}%
  \addtolength\belowdisplayshortskip{-0.5\baselineskip}%
}
%\usepackage[listformat=simple]{caption}
\usepackage{chngcntr}
\counterwithout{figure}{chapter}
\usepackage{pgfplots}
\pgfplotsset{compat=1.9}
% We will externalize the figures
\usepgfplotslibrary{external}
%\tikzexternalize

%%%%%%%%%%%%%%%%%%%%%%%%%%%%%%%%%%%%%%%%%%%%%%%%%%%%%%%%%%%%%%%%%%%%%%%%%%
%%%%%%%%%%%%%%%%%%%%%%%%%%%%%%%%%%%%%%%%%%%%%%%%%%%%%%%%%%%%%%%%%%%%%%%%%%
%%%%%%%%%%%%%%%%%%%%%%%%%%%%%%%%%%%%%%%%%%%%%%%%%%%%%%%%%%%%%%%%%%%%%%%%%%
%%%%%%%%%%%%%%%%%%%%%%%%%%%%%%%%%%%%%%%%%%%%%%%%%%%%%%%%%%%%%%%%%%%%%%%%%%
%
%  VOUS POUVEZ APPORTER QUELQUES AJOUTS ICI
%
%%%%%%%%%%%%%%%%%%%%%%%%%%%%%%%%%%%%%%%%%%%%%%%%%%%%%%%%%%%%%%%%%%%%%%%%%%
%%%%%%%%%%%%%%%%%%%%%%%%%%%%%%%%%%%%%%%%%%%%%%%%%%%%%%%%%%%%%%%%%%%%%%%%%%
%%%%%%%%%%%%%%%%%%%%%%%%%%%%%%%%%%%%%%%%%%%%%%%%%%%%%%%%%%%%%%%%%%%%%%%%%%
%%%%%%%%%%%%%%%%%%%%%%%%%%%%%%%%%%%%%%%%%%%%%%%%%%%%%%%%%%%%%%%%%%%%%%%%%%

%%%%%%%%%%%%%%%%%%%%%%%%%%%%%%%%%%%%%%%%%%%%%%%%%%%%%%%%%%%%%%%%%%%%%%%%%%
%%% MODIFIER ICI POUR VOTRE PDF SI BESOIN
%%%%%%%%%%%%%%%%%%%%%%%%%%%%%%%%%%%%%%%%%%%%%%%%%%%%%%%%%%%%%%%%%%%%%%%%%%
\usepackage[colorlinks=true]{hyperref}
\AtBeginDocument{%
 \hypersetup{
   urlcolor=linkEpisenMEMOIRE,
   urlbordercolor=linkEpisenMEMOIRE,
   linkcolor=black,
   linkbordercolor=linkEpisenMEMOIRE,
   citecolor=linkEpisenMEMOIRE,
   citebordercolor=linkEpisenMEMOIRE,
   anchorcolor=linkEpisenMEMOIRE,
   pdfborderstyle={/S/S/W 0.7},
   pdftitle=\TheManuscritTitle,
   pdfauthor=\TheAuthorManuscript\\ (\TheAuthorEmail),
   pdfsubject=\TheMemoireSubject,
   pdfkeywords=\TheMemoireKeywords
}}

%%%%%%%%%%%%%%%%%%%%%%%%%%%%%%%%%%%%%%%%%%%%%%%%%%%%%%%%%%%%%%%%%%%%%%%%%%
%%% MODIFIER ICI pour les package tikz
%%%%%%%%%%%%%%%%%%%%%%%%%%%%%%%%%%%%%%%%%%%%%%%%%%%%%%%%%%%%%%%%%%%%%%%%%%
\usetikzlibrary{automata,positioning}
\AtBeginEnvironment{tikzpicture}{\tracinglostchars=0\relax}   % ==> pour éviter des warnings sur un nullfont


%%%%%%%%%%%%%%%%%%%%%%%%%%%%%%%%%%%%%%%%%%%%%%%%%%%%%%%%%%%%%%%%%%%%%%%%%%
%  Ajoutez vos format de listes ici !
%%%%%%%%%%%%%%%%%%%%%%%%%%%%%%%%%%%%%%%%%%%%%%%%%%%%%%%%%%%%%%%%%%%%%%%%%%
\newenvironment{enumerateEPISEN}
{\begin{enumerate}[label=\color{colorEPISEN-MEMOIRE}\textbf{\arabic*.}]}
{\end{enumerate}}

\newenvironment{itemizeEPISEN}
{\begin{itemize}[label=\color{colorEPISEN-MEMOIRE}$\blacktriangleright$]
  \setlength{\itemsep}{0pt}%
  \setlength{\parskip}{0pt}}
{\end{itemize}}

\newenvironment{itemizeRondEPISEN}
{\begin{itemize}[label=\textbullet, font=\LARGE \color{colorEPISEN-MEMOIRE}] 
  \setlength{\itemsep}{0pt}%
  \setlength{\parskip}{0pt}}
{\end{itemize}}


%%%%%%%%%%%%%%%%%%%%%%%%%%%%%%%%%%%%%%%%%%%%%%%%%%%%%%%%%%%%%%%%%%%%%%%%%%
%%% Quelques commandes
%%%%%%%%%%%%%%%%%%%%%%%%%%%%%%%%%%%%%%%%%%%%%%%%%%%%%%%%%%%%%%%%%%%%%%%%%%

% ...
% ...
% ...
% ...


%%%%%%%%%%%%%%%%%%%%%%%%%%%%%%%%%%%%%%%%%%%%%%%%%%%%%%%%%%%%%%%%%%%%%%%%%%
%%% Programme
%%%%%%%%%%%%%%%%%%%%%%%%%%%%%%%%%%%%%%%%%%%%%%%%%%%%%%%%%%%%%%%%%%%%%%%%%%



%https://tex.stackexchange.com/questions/64900/improper-alphabetic-constant
% Lien très intéressant sur prendre directement un fichier et sélectionner
% que certaines lignes ... (TODO) Tester aussi pour des codes accessible sur GIT
% pour faire de l'open Science !!!!

%%% Java
%%%%%%%%%%%%%%%%%%%%%%%%%%%%%%%%%%%%%%%%%%%%%%%%%%%%%%%%%%%%%%%%%%%%%%%%%%
\lstdefinestyle{myJava}{language=java, 
	backgroundcolor=\color{linkEpisenMemoireUltraLight!30},
	keywordstyle=\bfseries\color{colorEPISEN-MEMOIRE},
	basicstyle=\normalfont,
	commentstyle=\color{linkEpisenGris},
	stringstyle=\color{brown},
    keepspaces=true, 
    numbers=left,
    numbersep=9pt,
    numberstyle=\tiny\color{linkEpisenGris},
    xleftmargin=20pt,
    framexleftmargin=1.77em,
    showstringspaces=true,
    showtabs=false,
    tabsize=2,
	breaklines=true,         % sets automatic line breaking
	breakatwhitespace=false, % sets if automatic breaks should only happen at whitespace
	columns=flexible,
	morekeywords={String, sealed, permits, System,out,println,-,+,*,/,++,--,<,>,=,!,!=},
	otherkeywords={\%*\%,+,!,!=,\&,<,>,=,++,--,\%/\%,\%\%},
	identifierstyle=\normalfont\itshape,
	literate=%
        {é}{{\'e}}1
        {è}{{\`e}}1
        {ê}{{\^e}}1
        {à}{{\`a}}1
        {â}{{\^a}}1
        {î}{{\^i}}1
        {ç}{{c}}1
        {ù}{{\`u}}1
    	{|}{{$\hskip .001em\smash{\color{linkEpisenGris}\Big|}\hskip .5em$}{ }}{1},
    mathescape=true,
}

\lstnewenvironment{java}
   {\lstset{style=myJava}}
   {}

%%% R
%%%%%%%%%%%%%%%%%%%%%%%%%%%%%%%%%%%%%%%%%%%%%%%%%%%%%%%%%%%%%%%%%%%%%%%%%%
\lstdefinelanguage{R}{
	keywords={abbreviate,abline,abs,acos,acosh,action,add1,add,%
      aggregate,alias,Alias,alist,all,anova,any,aov,aperm,append,apply,%
      approx,approxfun,apropos,Arg,args,array,arrows,as,asin,asinh,%
      atan,atan2,atanh,attach,attr,attributes,autoload,autoloader,ave,%
      axis,backsolve,barplot,basename,besselI,besselJ,besselK,besselY,%
      beta,binomial,body,box,boxplot,break,browser,bug,builtins,bxp,by,%
      c,C,call,Call,case,cat,category,cbind,ceiling,character,char,%
      charmatch,check,chol,chol2inv,choose,chull,class,close,cm,codes,%
      coef,coefficients,co,col,colnames,colors,colours,commandArgs,%
      comment,complete,complex,conflicts,Conj,contents,contour,%
      contrasts,contr,control,helmert,contrib,convolve,cooks,coords,%
      distance,coplot,cor,cos,cosh,count,fields,cov,covratio,wt,CRAN,%
      create,crossprod,cummax,cummin,cumprod,cumsum,curve,cut,cycle,D,%
      data,dataentry,date,dbeta,dbinom,dcauchy,dchisq,de,debug,%
      debugger,Defunct,default,delay,delete,deltat,demo,de,density,%
      deparse,dependencies,Deprecated,deriv,description,detach,%
      dev2bitmap,dev,cur,deviance,off,prev,,dexp,df,dfbetas,dffits,%
      dgamma,dgeom,dget,dhyper,diag,diff,digamma,dim,dimnames,dir,%
      dirname,dlnorm,dlogis,dnbinom,dnchisq,dnorm,do,dotplot,double,%
      download,dpois,dput,drop,drop1,dsignrank,dt,dummy,dump,dunif,%
      duplicated,dweibull,dwilcox,dyn,edit,eff,effects,eigen,else,%
      emacs,end,environment,env,erase,eval,equal,evalq,example,exists,%
      exit,exp,expand,expression,External,extract,extractAIC,factor,%
      fail,family,fft,file,filled,find,fitted,fivenum,fix,floor,for,%
      For,formals,format,formatC,formula,Fortran,forwardsolve,frame,%
      frequency,ftable,ftable2table,function,gamma,Gamma,gammaCody,%
      gaussian,gc,gcinfo,gctorture,get,getenv,geterrmessage,getOption,%
      getwd,gl,glm,globalenv,gnome,GNOME,graphics,gray,grep,grey,grid,%
      gsub,hasTsp,hat,heat,help,hist,home,hsv,httpclient,I,identify,if,%
      ifelse,Im,image,\%in\%,index,influence,measures,inherits,install,%
      installed,integer,interaction,interactive,Internal,intersect,%
      inverse,invisible,IQR,is,jitter,kappa,kronecker,labels,lapply,%
      layout,lbeta,lchoose,lcm,legend,length,levels,lgamma,library,%
      licence,license,lines,list,lm,load,local,locator,log,log10,log1p,%
      log2,logical,loglin,lower,lowess,ls,lsfit,lsf,ls,machine,Machine,%
      mad,mahalanobis,make,link,margin,match,Math,matlines,mat,matplot,%
      matpoints,matrix,max,mean,median,memory,menu,merge,methods,min,%
      missing,Mod,mode,model,response,mosaicplot,mtext,mvfft,na,nan,%
      names,omit,nargs,nchar,ncol,NCOL,new,next,NextMethod,nextn,%
      nlevels,nlm,noquote,NotYetImplemented,NotYetUsed,nrow,NROW,null,%
      numeric,\%o\%,objects,offset,old,on,Ops,optim,optimise,optimize,%
      options,or,order,ordered,outer,package,packages,page,pairlist,%
      pairs,palette,panel,par,parent,parse,paste,path,pbeta,pbinom,%
      pcauchy,pchisq,pentagamma,persp,pexp,pf,pgamma,pgeom,phyper,pico,%
      pictex,piechart,Platform,plnorm,plogis,plot,pmatch,pmax,pmin,%
      pnbinom,pnchisq,pnorm,points,poisson,poly,polygon,polyroot,pos,%
      postscript,power,ppoints,ppois,predict,preplot,pretty,Primitive,%
      print,prmatrix,proc,prod,profile,proj,prompt,prop,provide,%
      psignrank,ps,pt,ptukey,punif,pweibull,pwilcox,q,qbeta,qbinom,%
      qcauchy,qchisq,qexp,qf,qgamma,qgeom,qhyper,qlnorm,qlogis,qnbinom,%
      qnchisq,qnorm,qpois,qqline,qqnorm,qqplot,qr,Q,qty,qy,qsignrank,%
      qt,qtukey,quantile,quasi,quit,qunif,quote,qweibull,qwilcox,%
      rainbow,range,rank,rbeta,rbind,rbinom,rcauchy,rchisq,Re,read,csv,%
      csv2,fwf,readline,socket,real,Recall,rect,reformulate,regexpr,%
      relevel,remove,rep,repeat,replace,replications,report,require,%
      resid,residuals,restart,return,rev,rexp,rf,rgamma,rgb,rgeom,R,%
      rhyper,rle,rlnorm,rlogis,rm,rnbinom,RNGkind,rnorm,round,row,%
      rownames,rowsum,rpois,rsignrank,rstandard,rstudent,rt,rug,runif,%
      rweibull,rwilcox,sample,sapply,save,scale,scan,scan,screen,sd,se,%
      search,searchpaths,segments,seq,sequence,setdiff,setequal,set,%
      setwd,show,sign,signif,sin,single,sinh,sink,solve,sort,source,%
      spline,splinefun,split,sqrt,stars,start,stat,stem,step,stop,%
      storage,strstrheight,stripplot,strsplit,structure,strwidth,sub,%
      subset,substitute,substr,substring,sum,summary,sunflowerplot,svd,%
      sweep,switch,symbol,symbols,symnum,sys,status,system,t,table,%
      tabulate,tan,tanh,tapply,tempfile,terms,terrain,tetragamma,text,%
      time,title,topo,trace,traceback,transform,tri,trigamma,trunc,try,%
      ts,tsp,typeof,unclass,undebug,undoc,union,unique,uniroot,unix,%
      unlink,unlist,unname,untrace,update,upper,url,UseMethod,var,%
      variable,vector,Version,vi,warning,warnings,weighted,weights,%
      which,while,window,write,\%x\%,x11,X11,xedit,xemacs,xinch,xor,%
      xpdrows,xy,xyinch,yinch,zapsmall,zip},
	otherkeywords={0,1,2,3,4,5,6,7,8,9},
	morekeywords={TRUE,FALSE,for,if,then,else,function,in},
	otherkeywords={-,+,!,!=,<,>,~,$,*,\&,\%/\%,\%*\%,\%\%,_,/},%
	alsoother={._$},%
	sensitive,%
	morecomment=[l]{\#},%
	morestring=[d]",%
	morestring=[d]',
	literate=%
        {é}{{\'e}}1
        {è}{{\`e}}1
        {ê}{{\^e}}1
        {à}{{\`a}}1
        {â}{{\^a}}1
        {î}{{\^i}}1
        {ç}{{c}}1
        {ù}{{\`u}}1
    	{|}{{$\hskip .001em\smash{\color{linkEpisenGris}\Big|}\hskip .5em$}{ }}{1},
}

\lstdefinestyle{myR}{
	language=R,
	backgroundcolor=\color{linkEpisenMemoireUltraLight!30},
	keywordstyle=\bfseries\color{colorEPISEN-MEMOIRE},
	basicstyle=\normalfont,
	stringstyle=\color{brown},
	commentstyle=\color{linkEpisenGris},
	keepspaces=true, 
    numbers=left,
    numbersep=9pt,
    numberstyle=\tiny\color{linkEpisenGris},
    xleftmargin=20pt,
    framexleftmargin=1.5em,
    showstringspaces=true,
    tabsize=2,
	breaklines=true,         % sets automatic line breaking
	breakatwhitespace=false, % sets if automatic breaks should only happen at whitespace
	columns=flexible,
    escapeinside={\%*}{*)},
    identifierstyle=\normalfont\itshape,
} 

\lstnewenvironment{R}
    {\lstset{style=myR}}
    {}

%%% Python
%%%%%%%%%%%%%%%%%%%%%%%%%%%%%%%%%%%%%%%%%%%%%%%%%%%%%%%%%%%%%%%%%%%%%%%%%%

\lstdefinestyle{myPython}{language=python, 
	backgroundcolor=\color{linkEpisenMemoireUltraLight!30},
	keywordstyle=\bfseries\color{colorEPISEN-MEMOIRE},
	basicstyle=\normalfont,
	commentstyle=\color{linkEpisenGris},
	stringstyle=\color{brown},
    keepspaces=true, 
    numbers=left,
    numbersep=9pt,
    numberstyle=\tiny\color{linkEpisenGris},
    xleftmargin=20pt,
    framexleftmargin=1.5em,
    showstringspaces=true,
    showtabs=false,
    tabsize=2,
	breaklines=true,         % sets automatic line breaking
	breakatwhitespace=false, % sets if automatic breaks should only happen at whitespace
	columns=flexible,
	morekeywords={-,+,*,/,++,--,<,>,=,!,!=},
	otherkeywords={\%*\%,+,!,!=,\&,<,>,=,\%/\%,\%\%},
	identifierstyle=\normalfont\itshape,
	literate=%
        {é}{{\'e}}1
        {è}{{\`e}}1
        {ê}{{\^e}}1
        {à}{{\`a}}1
        {â}{{\^a}}1
        {î}{{\^i}}1
        {ç}{{c}}1
        {ù}{{\`u}}1
    	{|}{{$\hskip .001em\smash{\color{linkEpisenGris}\Big|}\hskip .5em$}{ }}{1},
    mathescape=true,
}

\lstnewenvironment{python}
   {\lstset{style=myPython}}
   {}


%%% Algorithmes
%%%%%%%%%%%%%%%%%%%%%%%%%%%%%%%%%%%%%%%%%%%%%%%%%%%%%%%%%%%%%%%%%%%%%%%%%%

%https://tex.stackexchange.com/questions/668747/different-color-for-algorithm2e-keywords
\renewcommand{\KwSty}[1]{\textnormal{\textcolor{colorEPISEN-MEMOIRE}{\bfseries #1}}\unskip}


% ****************
% Pour les autres langages, débrouillez vous un peu !!!
% ****************







%%%%%%%%%%%%%%%%%%%%%%%%%%%%%%%%%%%%%%%%%%%%%%%%%%%%%%%%%%%%%%%%%%%%%%%%%%
%%%%%%%%%%%%%%%%%%%%%%%%%%%%%%%%%%%%%%%%%%%%%%%%%%%%%%%%%%%%%%%%%%%%%%%%%%
%%%%%%%%%%%%%%%%%%%%%%%%%%%%%%%%%%%%%%%%%%%%%%%%%%%%%%%%%%%%%%%%%%%%%%%%%%
%%%%%%%%%%%%%%%%%%%%%%%%%%%%%%%%%%%%%%%%%%%%%%%%%%%%%%%%%%%%%%%%%%%%%%%%%%
%
%  A PRIORI, NE PLUS MODIFIER LA SUITE DE CE FICHIER !
%
%%%%%%%%%%%%%%%%%%%%%%%%%%%%%%%%%%%%%%%%%%%%%%%%%%%%%%%%%%%%%%%%%%%%%%%%%%
%%%%%%%%%%%%%%%%%%%%%%%%%%%%%%%%%%%%%%%%%%%%%%%%%%%%%%%%%%%%%%%%%%%%%%%%%%
%%%%%%%%%%%%%%%%%%%%%%%%%%%%%%%%%%%%%%%%%%%%%%%%%%%%%%%%%%%%%%%%%%%%%%%%%%
%%%%%%%%%%%%%%%%%%%%%%%%%%%%%%%%%%%%%%%%%%%%%%%%%%%%%%%%%%%%%%%%%%%%%%%%%%
\urlstyle{same}

%%%%%%%%%%%%%%%%%%%%%%%%%%%%%%%%%%%%%%%%%%%%%%%%%%%%%%%%%%%%%%%%%%%%%%%%%%
%%% Format Glossaire
%%%%%%%%%%%%%%%%%%%%%%%%%%%%%%%%%%%%%%%%%%%%%%%%%%%%%%%%%%%%%%%%%%%%%%%%%%
\usepackage[toc,acronym,style=listgroup]{glossaries}
% Autre style de glossaires ici : https://www.dickimaw-books.com/gallery/glossaries-styles/
\usepackage[automake]{glossaries-extra}
% ===> POUR OVERLEAF
%\usepackage[toc,acronym]{glossaries-extra}

%\newglossarystyle{dottedlocations}{%
%  \setglossarystyle{list}%
%  \renewcommand*{\glossentry}[2]{%
%    \item[\glsentryitem{##1}\glstarget{##1}{\glossentryname{##1}}] ##3%
%      \unskip\leaders\hbox to 2.9mm{\hss{\color{colorEPISEN-MEMOIRE}\LARGE.}}\hfill##5}
%  \renewcommand*{\glsgroupskip}{}%
%}

\newglossarystyle{dottedlocations}{%
    \setglossarystyle{list}%
    \renewcommand*{\glossentry}[2]{%
        \item[\glsentryitem{##1}\glstarget{##1}{\glossentryname{##1}}]
        \ifglshassymbol{##1}{[\glossentrysymbol{##1}]\quad}{}%
        {\glossentrydesc{##1}}%
        \unskip\leaders\hbox to 2.9mm{\hss{\color{colorEPISEN-MEMOIRE}\LARGE.}}\hfill##2}%
    \renewcommand*{\glsgroupskip}{}%
}
\setglossarystyle{dottedlocations}


\renewcommand*{\glsgroupheading}[1]{\item[]{\Huge\color{colorEPISEN-MEMOIRE}\textbf{\glsgetgrouptitle{#1}}}}

\renewcommand*{\glsnumberformat}[1]{\hyperlink{page.#1}{#1}}

\renewcommand{\glsnamefont}[1]{\color{colorEPISEN-MEMOIRE}\textbf{#1}}
\setlength{\glslistdottedwidth}{.3\linewidth}
\GlsXtrEnablePreLocationTag{ {\color{colorEPISEN-MEMOIRE}\textbf{Page~: }}}{{\color{colorEPISEN-MEMOIRE}\textbf{Pages~: }}}


%%%%%%%%%%%%%%%%%%%%%%%%%%%%%%%%%%%%%%%%%%%%%%%%%%%%%%%%%%%%%%%%%%%%%%%%%%
%%%%%%%%%%%%%%%%%%%%%%%%%%%%%%%%%%%%%%%%%%%%%%%%%%%%%%%%%%%%%%%%%%%%%%%%%%

\AtEndDocument{\label{LastPage}}

\linepenalty=1000
\everypar=\expandafter{\the\everypar\loosness=-1 }


%%%%%%%%%%%%%%%%%%%%%%%%%%%%%%%%%%%%%%%%%%%%%%%%%%%%%%%%%%%%%%%%%%%%%%%%%%
%%%%%%%%%%%%%%%%%%%%%%%%%%%%%%%%%%%%%%%%%%%%%%%%%%%%%%%%%%%%%%%%%%%%%%%%%%
%  BIBLIO
%%%%%%%%%%%%%%%%%%%%%%%%%%%%%%%%%%%%%%%%%%%%%%%%%%%%%%%%%%%%%%%%%%%%%%%%%%
%%%%%%%%%%%%%%%%%%%%%%%%%%%%%%%%%%%%%%%%%%%%%%%%%%%%%%%%%%%%%%%%%%%%%%%%%%
\ifthenelse{\equal{\TheLanguage}{french}}
 {
	\DefineBibliographyStrings{french}{%
  		backrefpage = {cf. page},
  		backrefpages = {cf. pages},
  		urlseen = {disponible~:},
  		urlfrom = {disponible à l'\textsc{url}~:},
  		in={dans},
  		volume={vol\adddot},
  		pages = {pp\adddot},
  		byeditor = {Sous la direction de}
 	}
 	\DefineBibliographyExtras{french}{\protected\def\bibrangedash{--}}
 }
 {
 	\DefineBibliographyStrings{english}{%
  		backrefpage = {see page},
  		backrefpages = {see pages},
  		urlseen = {Accessed:},
  		urlfrom = {Available at \textsc{url}:},
  		byeditor = {Edited by}
 	}
 }
\DeclareFieldFormat{url}{\bibstring{urlfrom}\addcolon\space\url{#1}}

\usepackage{xpatch}
% No dot before number of articles
\xpatchbibmacro{volume+number+eid}{%
  \setunit*{\adddot}%
}{%
}{}{}


% Comma before and after journal volume
\renewbibmacro*{volume+number+eid}{%
  \setunit*{\addcomma\space}% NEW
  \printfield{volume}%
%  \setunit*{\adddot}% DELETED
  \setunit*{\addcomma\space}% NEW
  \printfield{number}%
  \setunit{\addcomma\space}%
  \printfield{eid}}

% Prefixes for journal volume and number
\DeclareFieldFormat[article]{volume}{\bibstring{volume}~#1}% volume of a journal
\DeclareFieldFormat[article]{number}{\bibstring{number}~#1,}% number of a journal


\newtoggle{bbx:datemissing}

\renewbibmacro*{date}{\toggletrue{bbx:datemissing}%
}

\renewbibmacro*{issue+date}{%
  \toggletrue{bbx:datemissing}%
  \iffieldundef{issue}
    {}
    {\printtext[parens]{%
       \printfield{issue}}}%
  \newunit}

\newbibmacro*{date:print}{%
  \togglefalse{bbx:datemissing}%
  \printdate}

\renewbibmacro*{chapter+pages}{%
  \printfield{chapter}%
  \setunit{\bibpagespunct}%
  \printfield{pages}%
  \newunit
  \setunit*{\addcomma\space}
  \usebibmacro{date:print}%
  \newunit}

\renewbibmacro*{note+pages}{%
  \printfield{note}%
  \setunit{\bibpagespunct}%
  \printfield{pages}%
  \newunit
  \setunit*{\addcomma\space}
  \usebibmacro{date:print}%
  \newunit}

\renewbibmacro*{addendum+pubstate}{%
  \iftoggle{bbx:datemissing}
    {\usebibmacro{date:print}}
    {}%
  \printfield{addendum}%
  \newunit\newblock
  \printfield{pubstate}}


%%%%%%%%%%%%%%%%%%%%%%%%%%%%%%%%%%%%%%%%%%%%%%%%%%%%%%%%%%%%%%%%%%%%%%%%%%
%%%%%%%%%%%%%%%%%%%%%%%%%%%%%%%%%%%%%%%%%%%%%%%%%%%%%%%%%%%%%%%%%%%%%%%%%%
%  Lecture des fichiers et données partagées
%%%%%%%%%%%%%%%%%%%%%%%%%%%%%%%%%%%%%%%%%%%%%%%%%%%%%%%%%%%%%%%%%%%%%%%%%%
%%%%%%%%%%%%%%%%%%%%%%%%%%%%%%%%%%%%%%%%%%%%%%%%%%%%%%%%%%%%%%%%%%%%%%%%%%

\CatchFileDef\TheManuscritTitle{../titre.tex}{\endlinechar=-1}
\CatchFileDef\TheAuthorEmail{../email.tex}{\endlinechar=-1}
\CatchFileDef\TheAuthorManuscript{../auteur.tex}{\endlinechar=-1}
\CatchFileDef\TheMemoireSubject{../sujet.tex}{\endlinechar=-1}
\CatchFileDef\TheMemoireKeywords{../keywords.tex}{\endlinechar=-1}

\newcommand{\departementEPISEN}{SI}    % write SI or ITS or ISBS

%%%%%%%%%%%%%%%%%%%%%%%%%%%%%%%%%%%%%%%%%%%%%%%%%%%%%%%%%%%%%%%%%%%%%%%%%%
%%%%%%%%%%%%%%%%%%%%%%%%%%%%%%%%%%%%%%%%%%%%%%%%%%%%%%%%%%%%%%%%%%%%%%%%%%
%%%%%%%%%%%%%%%%%%%%%%%%%%%%%%%%%%%%%%%%%%%%%%%%%%%%%%%%%%%%%%%%%%%%%%%%%%
%%%%%%%%%%%%%%%%%%%%%%%%%%%%%%%%%%%%%%%%%%%%%%%%%%%%%%%%%%%%%%%%%%%%%%%%%%
%
%  A PRIORI, NE PAS MODIFIER CE FICHIER !
%
%%%%%%%%%%%%%%%%%%%%%%%%%%%%%%%%%%%%%%%%%%%%%%%%%%%%%%%%%%%%%%%%%%%%%%%%%%
%%%%%%%%%%%%%%%%%%%%%%%%%%%%%%%%%%%%%%%%%%%%%%%%%%%%%%%%%%%%%%%%%%%%%%%%%%
%%%%%%%%%%%%%%%%%%%%%%%%%%%%%%%%%%%%%%%%%%%%%%%%%%%%%%%%%%%%%%%%%%%%%%%%%%
%%%%%%%%%%%%%%%%%%%%%%%%%%%%%%%%%%%%%%%%%%%%%%%%%%%%%%%%%%%%%%%%%%%%%%%%%%


%%%%%%%%%%%%%%%%%%%%%%%%%%%%%%%%%%%%%%%%%%%%%%%%%%%%%%%%%%%%%%%%%%%%%%%%%%
%%%%%%%%%%%%%%%%%%%%%%%%%%%%%%%%%%%%%%%%%%%%%%%%%%%%%%%%%%%%%%%%%%%%%%%%%%
% *********************** LES COULEURS (Départements) ********************
%%%%%%%%%%%%%%%%%%%%%%%%%%%%%%%%%%%%%%%%%%%%%%%%%%%%%%%%%%%%%%%%%%%%%%%%%%
%%%%%%%%%%%%%%%%%%%%%%%%%%%%%%%%%%%%%%%%%%%%%%%%%%%%%%%%%%%%%%%%%%%%%%%%%%

\definecolor{linkEpisenGris}{RGB}{81 98 111}
\definecolor{colorEPISEN-UPEC}{HTML}{4c6172}
%\definecolor{cadmiumred}{rgb}{0.89, 0.0, 0.13}
\definecolor{redUPEC}{HTML}{ef293d}

\ifthenelse{\equal{\departementEPISEN}{ITS}}
  {
   %%%%%%%%%%%%%%%%%%%%%%%%%%% POUR ITS
   \definecolor{colorEPISEN-MEMOIRE}{HTML}{655aa0}
   \definecolor{linkEpisenMEMOIRE}{RGB}{101 90 160}
   \definecolor{linkEpisenPoster}{RGB}{101 90 160}
   %\definecolor{linkEpisenPosterLight}{RGB}{223 221 238}
   %\definecolor{linkEpisenPosterLight}{HTML}{dfddee}
   \colorlet{linkEpisenPosterLight}{linkEpisenPoster!20}
   \colorlet{linkEpisenMemoireUltraLight}{linkEpisenPoster!7}
  }
  {
   \ifthenelse{\equal{\departementEPISEN}{SI}}
    {
	%%%%%%%%%%%%%%%%%%%%%%%%%%% POUR SI
	\definecolor{colorEPISEN-MEMOIRE}{HTML}{009fe3}
	\definecolor{linkEpisenMEMOIRE}{RGB}{0 159 227}
	\definecolor{linkEpisenPoster}{RGB}{0 159 227}
	%\definecolor{linkEpisenPosterLight}{RGB}{204 239 252}
  %\definecolor{linkEpisenPosterLight}{HTML}{cceffc}
  \colorlet{linkEpisenPosterLight}{linkEpisenPoster!20}
  \colorlet{linkEpisenMemoireUltraLight}{linkEpisenPoster!7}
    }
    {
	%%%%%%%%%%%%%%%%%%%%%%%%%%% POUR ISBS
	\definecolor{colorEPISEN-MEMOIRE}{HTML}{2fac66}
	\definecolor{linkEpisenMEMOIRE}{RGB}{47 172 102}
	\definecolor{linkEpisenPoster}{RGB}{47 172 102}
  %\definecolor{linkEpisenPosterLight}{RGB}{217 241 225}
  %\definecolor{linkEpisenPosterLight}{HTML}{d9f1e1}
  \colorlet{linkEpisenPosterLight}{linkEpisenPoster!20}
  \colorlet{linkEpisenMemoireUltraLight}{linkEpisenPoster!7}
    }
  }


%\definecolor{linkEpisenCyan}{RGB}{0 159 227}
%\definecolor{linkEpisenVert}{RGB}{47 172 102}
%\definecolor{linkEpisenViolet}{RGB}{101 90 160}
%\definecolor{colorEPISEN-SI}{HTML}{009fe3}
%\definecolor{colorEPISEN-ISBS}{HTML}{2fac66}
%\definecolor{colorEPISEN-ITS}{HTML}{655aa0}
%\definecolor{colorEPISEN-UPEC}{HTML}{4c6172}


%%%%%%%%%%%%%%%%%%%%%%%%%%%%%%%%%%%%%%%%%%%%%%%%%%%%%%%%%%%%%%%%%%%%%%%%%%
%%%%%%%%%%%%%%%%%%%%%%%%%%%%%%%%%%%%%%%%%%%%%%%%%%%%%%%%%%%%%%%%%%%%%%%%%%
%  Format of the Chapter
%%%%%%%%%%%%%%%%%%%%%%%%%%%%%%%%%%%%%%%%%%%%%%%%%%%%%%%%%%%%%%%%%%%%%%%%%%
%%%%%%%%%%%%%%%%%%%%%%%%%%%%%%%%%%%%%%%%%%%%%%%%%%%%%%%%%%%%%%%%%%%%%%%%%%

\titleformat{\chapter}[display]
{\normalfont\color{colorEPISEN-MEMOIRE}}
{\filleft\huge\normalcolor\textsc\chaptertitlename\hspace*{2mm}%
	\begin{tikzpicture}[baseline={([yshift=-.6ex]current bounding box.center)}]
    \node[fill=colorEPISEN-MEMOIRE,circle,text=white] {\thechapter};
	\end{tikzpicture}
}
{1ex}
{\titlerule[2.5pt]\vspace*{1.5ex}\Huge\normalcolor\textsc}
[]


%%%%%%%%%%%%%%%%%%%%%%%%%%%%%%%%%%%%%%%%%%%%%%%%%%%%%%%%%%%%%%%%%%%%%%%%%%
%%%%%%%%%%%%%%%%%%%%%%%%%%%%%%%%%%%%%%%%%%%%%%%%%%%%%%%%%%%%%%%%%%%%%%%%%%
%  Format of the "Table of Contents"
%%%%%%%%%%%%%%%%%%%%%%%%%%%%%%%%%%%%%%%%%%%%%%%%%%%%%%%%%%%%%%%%%%%%%%%%%%
%%%%%%%%%%%%%%%%%%%%%%%%%%%%%%%%%%%%%%%%%%%%%%%%%%%%%%%%%%%%%%%%%%%%%%%%%%

% Chapter text styling
\titlecontents{chapter}
	[10.4mm] % Left indentation
	{\addvspace{10pt}\Large\sffamily\bfseries} % Spacing and font options for chapters
	{
	 \hspace*{-12.4mm}
	 \begin{tikzpicture}[baseline={([yshift=-.6ex]current bounding box.center)}]
       \node[fill=colorEPISEN-MEMOIRE,circle,text=white] {\thecontentslabel};
  	  \end{tikzpicture}
  	  %\color{colorEPISEN-MEMOIRE}\contentslabel[\Large\thecontentslabel]{5mm}\normalcolor} 
    }
	{} % Formatting of numberless sections of this type
	{\color{colorEPISEN-MEMOIRE}\normalsize\;\titlerule*[.5pc]{.}\;\Large\thecontentspage\normalcolor} % Formatting of the filler to the right of the heading and the page number


% Part text styling (this is mostly taken care of in the PART HEADINGS section of this file)
\titlecontents{part}
	[0cm] % Left indentation
	{\addvspace{20pt}\bfseries} % Spacing and font options for parts
	{}
	{}
	{}

% Section text styling
\titlecontents{section}
	[18.4mm] % Left indentation
	{\addvspace{4pt}\sffamily\bfseries} % Spacing and font options for sections
	{\contentslabel[\thecontentslabel]{7mm}} % Formatting of numbered sections of this type
	{} % Formatting of numberless sections of this type
	{\ \titlerule*[.5pc]{.}\;\thecontentspage} % Formatting of the filler to the right of the heading and the page number

% Subsection text styling
\titlecontents{subsection}
	[21.6mm] % Left indentation
	{\addvspace{3pt}\sffamily\small} % Spacing and font options for subsections
	{\contentslabel[\thecontentslabel]{7mm}} % Formatting of numbered sections of this type
	{} % Formatting of numberless sections of this type
	{\ \titlerule*[.5pc]{.}\;\thecontentspage} % Formatting of the filler to the right of the heading and the page number

\titleformat*{\section}{\LARGE\bfseries}
\titleformat*{\subsection}{\Large\bfseries}

\makeatletter
\renewcommand\@seccntformat[1]{\color{colorEPISEN-MEMOIRE} {\csname the#1\endcsname}\hspace{0.5em}}

 \renewcommand{\subsubsection}{\@startsection {subsubsection}{3}{\z@}
 {.2ex}
 {.2ex}
 {{\color{colorEPISEN-MEMOIRE} \hspace*{-18pt} $\blacktriangleright$\ } \normalfont\bfseries}}

\makeatother



%%%%%%%%%%%%%%%%%%%%%%%%%%%%%%%%%%%%%%%%%%%%%%%%%%%%%%%%%%%%%%%%%%%%%%%%%%
%%%%%%%%%%%%%%%%%%%%%%%%%%%%%%%%%%%%%%%%%%%%%%%%%%%%%%%%%%%%%%%%%%%%%%%%%%
%  Format of the "Table of Figures"
%%%%%%%%%%%%%%%%%%%%%%%%%%%%%%%%%%%%%%%%%%%%%%%%%%%%%%%%%%%%%%%%%%%%%%%%%%
%%%%%%%%%%%%%%%%%%%%%%%%%%%%%%%%%%%%%%%%%%%%%%%%%%%%%%%%%%%%%%%%%%%%%%%%%%

% Figure text styling
\titlecontents{figure}
	[15pt] % Left indentation
	{\vspace*{5pt}} % Spacing and font options for figures
	{
	\begin{tikzpicture}[baseline={([yshift=-.6ex]current bounding box.center)}]
       \node[fill=colorEPISEN-MEMOIRE,circle,text=white] {\thecontentslabel};
  	\end{tikzpicture}
	\hspace*{1pt}} % Formatting of numbered sections of this type
	{} % Formatting of numberless sections of this type
	{\color{colorEPISEN-MEMOIRE} \ \titlerule*[.5pc]{.}\;\normalcolor\thecontentspage} % Formatting of the filler to the right of the heading and the page number


\newcommand*{\noaddvspace}{\renewcommand*{\addvspace}[1]{}}
\addtocontents{lof}{\protect\noaddvspace}
\addtocontents{lof}{\linespread{1}\selectfont}




%%%%%%%%%%%%%%%%%%%%%%%%%%%%%%%%%%%%%%%%%%%%%%%%%%%%%%%%%%%%%%%%%%%%%%%%%%
%%%%%%%%%%%%%%%%%%%%%%%%%%%%%%%%%%%%%%%%%%%%%%%%%%%%%%%%%%%%%%%%%%%%%%%%%%
%  Format of the mini "Table of Contents"  (minitoc)
%%%%%%%%%%%%%%%%%%%%%%%%%%%%%%%%%%%%%%%%%%%%%%%%%%%%%%%%%%%%%%%%%%%%%%%%%%
%%%%%%%%%%%%%%%%%%%%%%%%%%%%%%%%%%%%%%%%%%%%%%%%%%%%%%%%%%%%%%%%%%%%%%%%%%

%command to print the actual minitoc
\newcommand{\printMyMinitoc}{%
	\noindent\hspace{1cm}%
	%\colorlet{chpnumbercolor}{black}%
	\begin{tikzpicture}
	\node(s){
	    \normalcolor
		\begin{minipage}{.9\linewidth}%minipage trick
		\begin{tocBox}		
		\printcontents[chapters]{}{1}{}
		\end{tocBox}
		\end{minipage}
	};
	{
        \color{colorEPISEN-MEMOIRE}

		\draw[line width=1.5pt](s.north west)--(s.north east) (s.south west)--(s.south east);
	}
	\end{tikzpicture}
}

\setlength\afterepigraphskip{-5pt}
\setlength\beforeepigraphskip{-20pt}

\let\originalepigraph\epigraph 
\renewcommand\epigraph[2]{\originalepigraph{\textit{#1}\color{colorEPISEN-MEMOIRE}}{\textsc{#2}}}


%%%%%%%%%%%%%%%%%%%%%%%%%%%%%%%%%%%%%%%%%%%%%%%%%%%%%%%%%%%%%%%%%%%%%%%%%%
%%%%%%%%%%%%%%%%%%%%%%%%%%%%%%%%%%%%%%%%%%%%%%%%%%%%%%%%%%%%%%%%%%%%%%%%%%
%  Format de la page de titre
%%%%%%%%%%%%%%%%%%%%%%%%%%%%%%%%%%%%%%%%%%%%%%%%%%%%%%%%%%%%%%%%%%%%%%%%%%
%%%%%%%%%%%%%%%%%%%%%%%%%%%%%%%%%%%%%%%%%%%%%%%%%%%%%%%%%%%%%%%%%%%%%%%%%%

\newcommand\titlepagedecoration{

	\begin{tikzpicture}[remember picture,overlay,shorten >= -10pt]
		
		\coordinate (aux1) at ([yshift=-284pt]current page.north east);
		\coordinate (aux2) at ([yshift=-679pt]current page.north east);
		\coordinate (aux3) at ([xshift=-4.5cm, yshift=-284pt]current page.north east) ;
		\coordinate (aux4) at ([yshift=-419pt]current page.north east);
		\coordinate (aux5) at ([yshift=-724pt]current page.north east);
		\coordinate (aux6) at ([yshift=-219pt]current page.north east);
		
		\begin{scope}[colorEPISEN-MEMOIRE,line width=12pt,rounded corners=12pt]
			\draw
			(aux1) -- coordinate (a)
			++(225:5) --
			++(-45:5.1) coordinate (b);
			\draw[shorten <= -10pt]
			(aux3) --
			(a) --
			(aux1);
			\draw[opacity=0.6,colorEPISEN-MEMOIRE!70,shorten <= -10pt]
			(b) --
			++(225:2.2) --
			++(-45:2.2);
		\end{scope}
		\draw[colorEPISEN-MEMOIRE!40,line width=8pt,rounded corners=8pt,shorten <= -10pt]
		(aux4) --
		++(225:0.8) --
		++(-45:0.8);
		\draw[colorEPISEN-MEMOIRE,line width=8pt,rounded corners=8pt,shorten <= -10pt]
		(aux5) --
		++(225:0.8) --
		++(-45:0.8);
		\draw[colorEPISEN-MEMOIRE,line width=8pt,rounded corners=8pt,shorten <= -10pt]
		(aux6) --
		++(225:0.8) --
		++(-45:0.8);
		\begin{scope}[colorEPISEN-MEMOIRE!70,line width=6pt,rounded corners=8pt]
			\draw[shorten <= -10pt, colorEPISEN-MEMOIRE!40]
			(aux2) --
			++(225:3) coordinate[pos=0.45] (c) --
			++(-45:3.1);
			\draw
			(aux2) --
			(c) --
			++(135:2.5) --
			++(45:2.5) --
			++(-45:2.5) coordinate[pos=0.3] (d);   
			\draw
			(d) -- +(45:1);
		\end{scope}
	\end{tikzpicture}%
}

%%%%%%%%%%%%%%%%%%%%%%%%%%%%%%%%%%%%%%%%%%%%%%%%%%%%%%%%%%%%%%%%%%%%%%%%%%
%%% Quelques commandes
%%%%%%%%%%%%%%%%%%%%%%%%%%%%%%%%%%%%%%%%%%%%%%%%%%%%%%%%%%%%%%%%%%%%%%%%%%

\newcommand*{\doi}[1]{\href{https://doi.org/\detokenize{#1}}{doi: \detokenize{#1}}}

\newcommand{\HRule}{\rule{\linewidth}{3.3mm}}
\newcommand{\Hrule}{\rule{\linewidth}{0.5mm}}

\newcommand{\cqfd}{$\blacksquare$}


%%%%%%%%%%%%%%%%%%%%%%%%%%%%%%%%%%%%%%%%%%%%%%%%%%%%%%%%%%%%%%%%%%%%%%%%%%
%%% les box
%%%%%%%%%%%%%%%%%%%%%%%%%%%%%%%%%%%%%%%%%%%%%%%%%%%%%%%%%%%%%%%%%%%%%%%%%%

\newenvironment{tocBox}
{\begin{mdframed}[skipabove=0pt,
		skipbelow=0pt,
		rightline=false,
		leftline=true,
		topline=false,
		bottomline=false,
		linecolor=colorEPISEN-MEMOIRE,
		backgroundcolor=linkEpisenMemoireUltraLight,
		innerleftmargin=5pt,
		innerrightmargin=5pt,
		innertopmargin=5pt,
		innerbottommargin=5pt,
		leftmargin=0cm,
		rightmargin=0cm,
		linewidth=5pt,
		align=right,
		userdefinedwidth=\textwidth
		]}
	{\end{mdframed}}


% Theorem box
\newmdenv[skipabove=7pt,
skipbelow=7pt,
backgroundcolor=black!5,
linecolor=colorEPISEN-MEMOIRE,
backgroundcolor=linkEpisenMemoireUltraLight,
innerleftmargin=5pt,
innerrightmargin=5pt,
innertopmargin=5pt,
leftmargin=0cm,
rightmargin=0cm,
innerbottommargin=5pt,
linewidth=2pt]{tBox}

% Definition box
\newmdenv[skipabove=7pt,
skipbelow=7pt,
rightline=false,
leftline=true,
topline=false,
bottomline=false,
backgroundcolor=gray!7,
linecolor=colorEPISEN-MEMOIRE,
innerleftmargin=5pt,
innerrightmargin=5pt,
innertopmargin=0pt,
leftmargin=0cm,
rightmargin=0cm,
linewidth=4pt,
innerbottommargin=0pt]{dBox}

% proposition box	  
\newmdenv[skipabove=7pt,
skipbelow=7pt,
rightline=false,
leftline=true,
topline=true,
bottomline=false,
backgroundcolor=gray!25,
linecolor=colorEPISEN-MEMOIRE,
innerleftmargin=5pt,
innerrightmargin=5pt,
innertopmargin=2pt,
innerbottommargin=5pt,
leftmargin=0cm,
rightmargin=0cm,
linewidth=2pt]{pBox}

% Lemma box	  
\newmdenv[skipabove=7pt,
skipbelow=7pt,
rightline=false,
leftline=true,
topline=false,
bottomline=false,
backgroundcolor=gray!15,
linecolor=colorEPISEN-MEMOIRE,
innerleftmargin=5pt,
innerrightmargin=5pt,
innertopmargin=0pt,
innerbottommargin=0pt,
leftmargin=0cm,
rightmargin=0cm,
linewidth=4pt]{lBox}

% Proof box
\newmdenv[skipabove=7pt,
skipbelow=7pt,
rightline=false,
leftline=false,
topline=false,
bottomline=false,
innerleftmargin=5pt,
innerrightmargin=5pt,
innertopmargin=0pt,
leftmargin=0cm,
rightmargin=0cm,
linewidth=4pt,
innerbottommargin=0pt]{proofBox}


% Exemple box
\newmdenv[skipabove=7pt,
skipbelow=7pt,
rightline=false,
leftline=true,
topline=false,
bottomline=false,
linecolor=colorEPISEN-MEMOIRE,
backgroundcolor=black!5,
innerleftmargin=5pt,
innerrightmargin=0pt,
innertopmargin=0pt,
leftmargin=0cm,
rightmargin=0cm,
linewidth=4pt,
innerbottommargin=0pt]{exampleBox}

% Code box
\newmdenv[skipabove=7pt,
skipbelow=7pt,
rightline=false,
leftline=true,
topline=false,
bottomline=false,
linecolor=colorEPISEN-MEMOIRE,
innerleftmargin=0pt,
innerrightmargin=0pt,
innertopmargin=0pt,
leftmargin=0cm,
rightmargin=0cm,
linewidth=4pt,
innerbottommargin=0pt]{codeBox}

%%%%%%%%%%%%%%%%%%%%%%%%%%%%%%%%%%%%%%%%%%%%%%%%%%%%%%%%%%%%%%%%%%%%%%%%%%
%%% les box pour les théoremes/lemma/exemples/codes
%%%%%%%%%%%%%%%%%%%%%%%%%%%%%%%%%%%%%%%%%%%%%%%%%%%%%%%%%%%%%%%%%%%%%%%%%%

\newtheoremstyle{blacknumex}% Theorem style name
{9pt}% Space above
{5pt}% Space below
{\normalfont}% Body font
{} % Indent amount
{\bf\sffamily}% Theorem head font
{\;}% Punctuation after theorem head
{0.25em}% Space after theorem head
{{\color{colorEPISEN-MEMOIRE}\thmname{#1}\thmnumber{ #2}}\thmnote{\bfseries\ ---\nobreakspace\textit{#3}}}

\theoremstyle{blacknumex}

%%%%%%%%%% Théoreme
\newtheorem{theoremeT}{Théorème}[part]
\newenvironment{theorem}{\begin{tBox}\begin{theoremeT}}{\end{theoremeT}\end{tBox}}
\renewcommand{\thetheoremeT}{\arabic{theoremeT}}

%%%%%%%%%% Définition
\newtheorem{definitionT}{Définition}[part]
\newenvironment{definition}{\begin{dBox}\begin{definitionT}}{\end{definitionT}\end{dBox}}
\renewcommand{\thedefinitionT}{\arabic{definitionT}}

%%%%%%%%%% Lemma
\newtheorem{lemmaT}{Lemme}[part]
\newenvironment{lemma}{\begin{lBox}\begin{lemmaT}}{\end{lemmaT}\end{lBox}}
\renewcommand{\thelemmaT}{\arabic{lemmaT}}


%%%%%%%%%% Proposition
\newtheorem{propositionT}{Proposition}[part]
\newenvironment{proposition}{\begin{pBox}\begin{propositionT}}{\end{propositionT}\end{pBox}}
\renewcommand{\thepropositionT}{\arabic{propositionT}}

%%%%%%%%%% Corolaire
\newtheorem{corollaryT}{Corolaire}[part]
\newenvironment{corollary}{\begin{pBox}\begin{corollaryT}}{\end{corollaryT}\end{pBox}}
\renewcommand{\thecorollaryT}{\arabic{corollaryT}}

\newtheoremstyle{proofStyle}% Theorem style name
{9pt}% Space above
{5pt}% Space below
{\normalfont}% Body font
{-6pt} % Indent amount
{\bf\sffamily}% Theorem head font
{\;}% Punctuation after theorem head
{0.25em}% Space after theorem head
{{\color{colorEPISEN-MEMOIRE}\thmname{#1}} {\bfseries\ ---}}

\theoremstyle{proofStyle}

%%%%%%%%%% Proof
% \def\myhfillEPISEN{
%   \parfillskip=0pt
%   \widowpenalty=10000
%   \displaywidowpenalty=10000
%   \finalhyphendemerits=0
%   \unskip\nobreak\null\hfil\penalty50
%   \hskip2em\null\hfill
% }
% \def\qedEPISEN{\myhfillEPISEN{\color{colorEPISEN-MEMOIRE}\cqfd}\par}

\newtheorem{proofT}{Preuve}[part]
\renewenvironment{proof}
 {\begin{proofBox}\begin{proofT}}
 {\hfill{\color{colorEPISEN-MEMOIRE}\cqfd}\end{proofT}\end{proofBox}}
\renewcommand{\theproofT}{\arabic{proofT}}

\newtheoremstyle{exampleStyle}
{9pt}% Space above
{5pt}% Space below
{\normalfont}% Body font
{-6pt} % Indent amount
{\bf\sffamily}% Theorem head font
{\\}% Punctuation after theorem head
{0.25em}% Space after theorem head
{{\color{colorEPISEN-MEMOIRE}\thmname{#1}}}

\theoremstyle{exampleStyle}

\newtheorem*{exampleS}{$\ $ Example}
\newenvironment{example}{\begin{exampleBox}\begin{exampleS}}{\end{exampleS}\end{exampleBox}}


\newtheoremstyle{emptyStyle}
{9pt}% Space above
{5pt}% Space below
{\normalfont}% Body font
{-6pt} % Indent amount
{\bf\sffamily}% Theorem head font
{ }% Punctuation after theorem head
{0.25em}% Space after theorem head
{}
\theoremstyle{emptyStyle}

\newtheorem*{codeS}{}
\newenvironment{code}{\begin{codeBox}\begin{codeS}}{\end{codeS}\end{codeBox}}
\newenvironment{codeExample}{\begin{exampleBox}\begin{codeS}}{\end{codeS}\end{exampleBox}}


%%% Mettre le glossaire !
%%%%%%%%%%%%%%%%%%%%%%%%%%%%%%%%%%%%%%%%%%%%%%%%%%%%%%%%%%%%%%%%%%%%%%%%%%
\makeglossaries

%******************************************************************************
%******************************** GLOSSAIRE ***********************************
%******************************************************************************

\newglossaryentry{prefixe}
{
    name=préfixe,
    sort=prefixe,
    description={Si le mot à indexé a un accent, alors definir son ordre sans accent. Et on tappe le mot sans accent. C'est lourd, mais je n'ai pas trouvé mieux...}
}


\newglossaryentry{latex}
{
    name=latex,
    description={Is a markup language specially suited 
    for scientific documents; PUIS ON AJOUTE DU BLABLA~; Quisque mi lorem, pulvinar eget, egestas quis, luctus at, ante. Proin auctor vehicula purus. Fusce ac nisl aliquam ante hendrerit pellentesque. Class aptent taciti sociosqu ad litora torquent per conubia nostra, per inceptos hymenaeos. Morbi wisi. Etiam arcu mauris, facilisis sed, eleifend non, nonummy ut, pede. Crasut lacus tempor metus mollis placerat. Vivamus eu tortor vel metus interdum malesuada.}
}

\newglossaryentry{EPISEN}
{
    name=EPISEN,
    description={Une école sympa, je pense que vous êtes d'accord avec cela ?; PUIS ON AJOUTE DU BLABLA~; Quisque mi lorem, pulvinar eget, egestas quis, luctus at, ante. Proin auctor vehicula purus. Fusce ac nisl aliquam ante hendrerit pellentesque. Class aptent taciti sociosqu ad litora torquent per conubia nostra, per inceptos hymenaeos. Morbi wisi. Etiam arcu mauris, facilisis sed, eleifend non, nonummy ut, pede. Crasut lacus tempor metus mollis placerat. Vivamus eu tortor vel metus interdum malesuada.}
}

%******************************************************************************
%******************************** Acronymes ***********************************
%******************************************************************************

\newacronym{gcd}{GCD}{{\bf\color{colorEPISEN-MEMOIRE} \textsc{g}}reatest {\bf\color{colorEPISEN-MEMOIRE} \textsc{c}}ommon {\bf\color{colorEPISEN-MEMOIRE} \textsc{d}}ivisor~; {\bf Rappel~:} de maths pour ceux au fond de salle que blablabla~; PUIS ON AJOUTE DU BLABLA  Quisque mi lorem, pulvinar eget, egestas quis, luctus at, ante. Proin auctor vehicula purus. Fusce ac nisl aliquam ante hendrerit pellentesque. Class aptent taciti sociosqu ad litora torquent per conubia nostra, per inceptos hymenaeos. Morbi wisi. Etiam arcu mauris, facilisis sed, eleifend non, nonummy ut, pede. Crasut lacus tempor metus mollis placerat. Vivamus eu tortor vel metus interdum malesuada.}

\newacronym{llm}{LLM}{{\bf\color{colorEPISEN-MEMOIRE} \textsc{l}}arge {\bf\color{colorEPISEN-MEMOIRE} \textsc{l}}anguage {\bf\color{colorEPISEN-MEMOIRE} \textsc{m}}odel~; {\bf Rappel~:} pour ceux qui ne lisent aucun journal que blablabla~; PUIS ON AJOUTE DU BLABLA  Duis nisl nibh, laoreet suscipit, convallis ut, rutrum id, enim. Phasellus odio. Nulla nulla elit, molestie non, scelerisque at, vestibulum eu, nulla. Ut odio nisl, facilisis id, mollis et, scelerisque nec, enim. Aenean sem leo, pellentesque sit amet, scelerisque sit amet, vehicula pellentesque, sapien.}

